\begin{definition}[The round's graded-agreement instance]
  \label{def:aba-gbca-instance-by-abdy}
  \lean{PLTS.ABA.GBCA.ByABDY.GBCAEvent}\lean{PLTS.ABA.GBCA.ByABDY.GBCALabel}\lean{PLTS.ABA.GBCA.ByABDY.gbcaEvents}\lean{PLTS.ABA.GBCA.ByABDY.GBCAProgramStep}\lean{PLTS.ABA.GBCA.ByABDY.gbcaProgram}\lean{PLTS.ABA.GBCA.ByABDY.NetworkState}\lean{PLTS.ABA.GBCA.ByABDY.NetworkState.corrupt}\lean{PLTS.ABA.GBCA.ByABDY.GBCANetworkStep}\lean{PLTS.ABA.GBCA.ByABDY.GBCANetwork}\lean{PLTS.ABA.GBCA.ByABDY.RoundState}\lean{PLTS.ABA.GBCA.ByABDY.compositionExtended}\lean{PLTS.ABA.GBCA.ByABDY.composition}\lean{PLTS.ABA.GBCA.ByABDY.roundOwnsLabel}\lean{PLTS.ABA.GBCA.ByABDY.corruptionAct}\lean{PLTS.ABA.GBCA.ByABDY.gbcaInstanceFamily}\leanok
  \uses{def:aba-programs, def:aba-network, def:aba-round-messages, def:relabel, def:synchronised-product, def:idle-family, def:process-algebra}
  One round of graded agreement taken apart into the pieces that run it: $n$ local graded-agreement programs $\mathrm{gbcaProgram}(P,r,j)$, each holding one process's round variables of Definition~\ref{def:aba-programs}, beside that round's own network $\mathrm{GBCANetwork}(P,r)$, which holds the round's per-sender message sets (D5), the corrupted set and the round's bound bit.
  That bit is a ghost output (D29): the round's first return writes it, as $\mathrm{boundOf}$ of the round's sent sets, the corrupted set and the outcome, a later return leaves it, and the network announces it on $\mathrm{boundG}(r,b)$ by a self-loop that may repeat; no program's variables hold it.
  The instance
  \[
    \mathrm{composition}(P,r) \;=\; \mathrm{relabel}\Bigl(
      \mathrm{abstract}_{\mathrm{GBCAEvent}(n)}\bigl(
        \bigl(\mathrm{synchronisedProduct}_{j \in [n]}\ \mathrm{gbcaProgram}(P,r,j)\bigr)
          \parallel \mathrm{GBCANetwork}(P,r)\bigr)\Bigr)
  \]
  is read over $\mathrm{GBCALabel}(n) = \mathrm{ExtendedLabel}(n, \mathrm{Message}) \oplus \mathrm{GBCAEvent}(n)$ and then over $\mathrm{ExtendedLabel}(n, \mathrm{Message})$, where $\mathrm{GBCAEvent}(n)$ is the alphabet of the two labels the components of one round synchronise on, the multicast and the delivery.
  They carry no round tag, the round being the identity of the instance they belong to, and they are hidden before they leave the instance, so its interface is the round's labels alone: $\mathrm{callG}(r,\cdot,\cdot)$, $\mathrm{retG}(r,\cdot,\cdot)$, $\mathrm{boundG}(r,\cdot)$ and a corrupted process's call $\mathrm{byzantineCallG}(r,\cdot,\cdot)$, whose effect the instance carries and whose guards $k \in F$ and the started rounds are read outside it (D11, D37).
  A program takes the call and a corrupted process's call alike, recording the input at a program holding none and looping at one holding an input.
  The returns and the announcement are the instance's outputs, and the calls its inputs, which it takes at every state.
  The graded-agreement component of the protocol is the $\mathbb{N}$-indexed family $\mathrm{gbcaInstanceFamily}(P)$ over $\mathrm{composition}(P,\cdot)$, in which the ownership map $\mathrm{roundOwnsLabel}$ sends each of those labels to its round, a round-tagged label moves its round alone, and $\mathrm{fail}(k)$ is the broadcast that keeps every round's copy of the corrupted set equal (D1).
  Its state holds the recorded corruptions and a finite map of the rounds the family has touched, a round enters the map at its first owned transition, and a round outside the map reads as the initial instance under the recorded corruptions (Definition~\ref{def:idle-family}).
  Full construction: \leandecl{PLTS.ABA.GBCA.ByABDY.composition} and \leandecl{PLTS.ABA.GBCA.ByABDY.gbcaInstanceFamily}.
\end{definition}
